\chapter{Worked Examples}\label{appA:examples}

This appendix computes, entirely by hand, three complete examples of the
kernel's machinery: an endomap on a three-element carrier taken through
its functional graph and its powers to idempotence; the deterministic
lift of the same map with all of its derived operators and its exact
norm; and a full closure run of the normal-form catalog on a two-element
carrier, with the shortlex table shown at every step. Nothing here
requires the software --- every line is derivable from the definitions
of Chapters~3--8 --- and that is the point: the examples are small
enough to audit the implementation against, and the bounded-exhaustive
suites of Chapter~10 are, in essence, these computations repeated over
every small instance.

\section{A three-carrier endomap and its powers}\label{appA:sec:endomap}

\subsection{The object and the map}

Fix the object $X = (\texttt{"X"}, (a, b, c))$: name \texttt{"X"},
ordered carrier $(a, b, c)$, standard indices $0, 1, 2$. Define the
endomap $f : X \to X$ by the mapping tuple

\[
f = (1, 2, 2), \qquad\text{i.e.}\qquad
f(a) = b,\quad f(b) = c,\quad f(c) = c .
\]

Construction succeeds: the name is non-blank, the carrier elements are
distinct and hashable, the mapping has length $3 = |X|$, and every image
index lies in $\{0, 1, 2\}$. Since source and target are the same
object, \texttt{is\_endomap} holds and the \texttt{Endomap} wrapper
accepts $f$.

\subsection{The functional graph}

The functional graph of $f$ has an edge $x \to f(x)$ from each index:

\[
0 \longrightarrow 1 \longrightarrow 2 \circlearrowleft
\]

Index $2$ is a fixed point --- a cycle of length $\ell = 1$ --- and
indices $0$ and $1$ form a tail feeding it. Tail lengths (distance to
the cycle): $d(0) = 2$, $d(1) = 1$, $d(2) = 0$. The maximum tail length
is $2 = |X| - 1$, the extreme the safe-power proof allows.

\subsection{Iterates}

Compute the iterates by composing mapping tuples,
$(g \circ f)(i) = g(f(i))$:

\begin{align*}
f^1 &= (1, 2, 2) \\
f^2 &= (f(1), f(2), f(2)) = (2, 2, 2) \\
f^3 &= (f(2), f(2), f(2)) = (2, 2, 2) = f^2 .
\end{align*}

So the iterates stabilize at the constant map to $c$ from exponent $2$
onward. Check idempotence of $f^2$ directly:
$f^2 \circ f^2 = (2,2,2) \circ (2,2,2)$; entry $i$ is
$f^2(f^2(i)) = f^2(2) = 2$, so $f^2 \circ f^2 = (2,2,2) = f^2$.
Idempotent. Note $f$ itself is \emph{not} idempotent:
$f \circ f = (2,2,2) \ne (1,2,2) = f$.

\subsection{The factorial power by binary powering}

The safe exponent is $|X|! = 3! = 6$. The implementation computes
$f^6$ by binary powering (Chapter~9); we trace it exactly. Write the
exponent in binary: $6 = 110_2$. Initialize the accumulator
$r = \id_X = (0, 1, 2)$ and the base $b = f = (1, 2, 2)$, with
$\mathrm{remaining} = 6$.

\begin{center}
\begin{tabular}{@{}clll@{}}
\toprule
Step & remaining (binary) & low bit & action \\
\midrule
1 & $110$ & $0$ & no accumulation; square: $b \leftarrow b \circ b = (2,2,2)$ \\
2 & $11$  & $1$ & $r \leftarrow b \circ r = (2,2,2)$; square: $b \leftarrow (2,2,2)\circ(2,2,2) = (2,2,2)$ \\
3 & $1$   & $1$ & $r \leftarrow b \circ r = (2,2,2)$; remaining becomes $0$, final squaring skipped \\
\bottomrule
\end{tabular}
\end{center}

In step 2, $b \circ r$ with $r = \id$ gives $b = (2,2,2)$; in step 3,
$(2,2,2) \circ (2,2,2) = (2,2,2)$ as computed above. The result is

\[
f^{6} = (2, 2, 2),
\]

which agrees with $f^2 = f^3 = \cdots$ from the iterate table, and is
idempotent as the safe-power theorem requires. The returned arrow
carries the presentation name \texttt{f\^{}(6)} --- the exponent is
well under the $256$-bit threshold, so the compact label is plain
decimal --- and, by extensional equality, equals the arrow $(2,2,2)$
regardless of its name.

\subsection{A contrasting endomap: a genuine cycle}

For contrast, let $g = (1, 0, 2)$: the transposition of $a, b$ with
$c$ fixed. Its functional graph is a $2$-cycle plus a fixed point;
there are no tails. Iterates: $g^2 = (g(1), g(0), g(2)) = (0, 1, 2) =
\id$, so $g^3 = g$, $g^4 = \id$, and in general the iterates alternate.
The safe power is $g^{6} = (g^2)^3 = \id^3 = \id$, which is idempotent
--- and here the factorial exponent is genuinely needed in the sense
that odd exponents never yield an idempotent: $g^n = g$ for odd $n$,
and $g \circ g = \id \ne g$. The two examples bracket the theorem's
content: for tail-dominated maps the iterates stabilize early and
$|X|!$ is generous; for permutation-like maps stabilization never
happens, but $|X|!$ is divisible by every cycle length ($1$ and $2$
here, with $2 \mid 6$) and so lands every point back at the start of
its cycle simultaneously.

\section{The lift of $f$, its operators, and its norm}\label{appA:sec:lift}

\subsection{The matrix}

Lift the endomap $f = (1,2,2)$ of the previous section. The matrix
$A_f$ is $|X| \times |X| = 3 \times 3$, with entry $(r, c)$ equal to
$1$ iff $f(c) = r$; column $c$ has its single $1$ in row $f(c)$:

\[
A_f \;=\;
\begin{pmatrix}
0 & 0 & 0 \\
1 & 0 & 0 \\
0 & 1 & 1
\end{pmatrix}
\]

(columns $0,1,2$ carry their $1$ in rows $1, 2, 2$ respectively).

\subsection{Fibers}

The fibers of $f$ are $f^{-1}(0) = \emptyset$, $f^{-1}(1) = \{0\}$,
$f^{-1}(2) = \{1, 2\}$, so

\[
\texttt{fiber\_sizes} = (0, 1, 2), \qquad
\texttt{max\_fiber\_size} = 2 .
\]

\subsection{The three derived operators}

\emph{Input Gram matrix} $A_f^{*} A_f$: entry $(x, x')$ is $1$ iff
$f(x) = f(x')$. The image tuple is $(1, 2, 2)$, so indices $1$ and $2$
collide and index $0$ is alone:

\[
A_f^{*} A_f \;=\;
\begin{pmatrix}
1 & 0 & 0 \\
0 & 1 & 1 \\
0 & 1 & 1
\end{pmatrix}
\]

\emph{Output Gram matrix} $A_f A_f^{*}$: diagonal of fiber sizes,

\[
A_f A_f^{*} \;=\; \mathrm{diag}(0, 1, 2) .
\]

\emph{Support projector}: diagonal $0/1$ marking of nonempty fibers,

\[
P_{\mathrm{supp}} \;=\; \mathrm{diag}(0, 1, 1) .
\]

Note $A_f A_f^{*} \ne P_{\mathrm{supp}}$ precisely because the fiber
over index $2$ has size two --- this instance witnesses the
three-operator distinction that Chapter~9 documents and the test suite
enforces.

\subsection{Scalar invariants}

Reading off the implemented consequences:

\[
\trace(A_f^{*} A_f) = 1 + 1 + 1 = 3 = |X|, \qquad
\rk(A_f) = |\{1, 2\}| = 2 = |\im f| ,
\]

with the trace computable either as the Gram diagonal sum (three ones)
or as the sum of fiber sizes $0 + 1 + 2 = 3$; both agree, as they
must. $f$ is not injective (fiber size $2$ over index $2$), so
\texttt{is\_injective}, \texttt{is\_isometry}, and
\texttt{is\_partial\_isometry} are all false. $f$ is an endomap but
$\texttt{mapping}[1] = 2 \ne 1$ for the image value $1$, so
\texttt{is\_idempotent} is false --- consistent with
$f \circ f \ne f$ above --- and \texttt{is\_orthogonal\_projector} is
false since the mapping is not the identity tuple.

\subsection{The exact norm, spectrally and by witness}

\emph{Spectral route.} Group the input basis by fibers: the singleton
fiber $\{0\}$ contributes the block $J_1 = (1)$ with eigenvalue $1$;
the fiber $\{1, 2\}$ contributes the all-ones block
$J_2 = \bigl(\begin{smallmatrix}1&1\\1&1\end{smallmatrix}\bigr)$
with eigenvalues $2$ and $0$. (Here the basis is already
fiber-grouped, so no permutation is needed.) The spectrum of
$A_f^{*} A_f$ is $\{1, 2, 0\}$, whence

\[
\lVert A_f \rVert^2 = \max\{1, 2, 0\} = 2
= \max_y |f^{-1}(y)|, \qquad
\lVert A_f \rVert = \sqrt{2} .
\]

\emph{Witness route.} The unit vector
$v = \bigl(0, \tfrac{1}{\sqrt 2}, \tfrac{1}{\sqrt 2}\bigr)$,
supported on the largest fiber, maps under the scatter-add rule to

\[
A_f v = \Bigl(0,\; 0,\; \tfrac{1}{\sqrt2} + \tfrac{1}{\sqrt2}\Bigr)
      = (0, 0, \sqrt 2),
\]

of norm $\sqrt 2$: the two colliding coordinates add coherently in the
shared output slot. So $\lVert A_f\rVert \ge \sqrt 2$, and the Gram
spectrum shows equality. The adjoint, for comparison, acts by
gathering: $A_f^{*}(w_0, w_1, w_2) = (w_1, w_2, w_2)$.

\subsection{An idempotent lift of norm $\sqrt 3$}

Finally, lift $h = f^2 = (2,2,2)$, the constant map to $c$. Its matrix
has a single nonzero row:

\[
A_h \;=\;
\begin{pmatrix}
0 & 0 & 0 \\
0 & 0 & 0 \\
1 & 1 & 1
\end{pmatrix},
\qquad
\texttt{fiber\_sizes} = (0, 0, 3) .
\]

Every image value is $2$ and $\texttt{mapping}[2] = 2$, so
\texttt{is\_idempotent} is \emph{true} --- indeed $A_h^2 = A_h$, since
row $2$ of $A_h A_h$ is again all ones. Yet
$\lVert A_h \rVert^2 = 3$, $\rk(A_h) = 1$,
$\trace(A_h^* A_h) = 3$, and $A_h$ is not symmetric, hence not an
orthogonal projector, and \texttt{is\_orthogonal\_projector} is false.
This is the two-state noncontractivity fixture of Chapter~7 scaled to
three states: an idempotent deterministic lift of norm $\sqrt 3 > 1$,
maximally far from a contraction, produced by two applications of a
perfectly ordinary endomap.

\section{A closure run with the shortlex table}\label{appA:sec:closure}

\subsection{The catalog}

Take one object $B = (\texttt{"B"}, (0, 1))$ and two generators, in
this declared order:

\begin{center}
\begin{tabular}{@{}llll@{}}
\toprule
Rank & Label & Mapping & Denotation \\
\midrule
$0$ & \texttt{s} & $(1, 0)$ & the swap \\
$1$ & \texttt{k} & $(0, 0)$ & the constant to carrier index $0$ \\
\bottomrule
\end{tabular}
\end{center}

Both generators are endo on $B$, so every extension is well-typed and
the closure explores freely. The four total maps on a two-element
carrier are $\id = (0,1)$, the swap $(1,0)$, and the two constants
$(0,0)$ and $(1,1)$; a complete run must find all four (this catalog
generates the full transformation monoid) and assign each its
shortlex-least word. Recall the word convention: a word's labels apply
left to right, so the word $(\texttt{k}, \texttt{s})$ denotes
$s \circ k$ --- first \texttt{k}, then \texttt{s}.

\subsection{The run, step by step}

\emph{Seeding.} The empty word $\varepsilon_B$ at $B$, denoting
$\id = (0,1)$, is pushed with key $(0, (), 0)$ (length $0$, empty rank
sequence, object rank $0$). Frontier: $\{\varepsilon_B\}$; queued
denotations: $\{(0,1)\}$.

\emph{Step 1: pop $\varepsilon_B$.} Its denotation $(0,1)$ is new;
accept it with representative $\varepsilon_B$. Extend by each
generator in rank order:
\begin{itemize}
\item by \texttt{s}: word $(\texttt{s})$, denotation
      $s \circ \id = (1,0)$ --- new; push with key $(1, (0), 0)$.
\item by \texttt{k}: word $(\texttt{k})$, denotation $(0,0)$ --- new;
      push with key $(1, (1), 0)$.
\end{itemize}
Table so far: $(0,1) \mapsto \varepsilon_B$. Queued: $(1,0), (0,0)$.

\emph{Step 2: pop $(\texttt{s})$} (key $(1,(0),0)$, the least in the
frontier). Accept $(1,0) \mapsto (\texttt{s})$. Extend:
\begin{itemize}
\item by \texttt{s}: denotation $s \circ s = (0,1)$ --- already
      represented (it is $\id$); skip, do not push.
\item by \texttt{k}: denotation $k \circ s = (k(1), k(0)) = (0,0)$
      --- not represented, but already \emph{queued} (the word
      $(\texttt{k})$ carries it); skip. This is the frontier
      deduplication of Chapter~9 doing its work: the queued word
      $(\texttt{k})$ has key $(1,(1),0)$, shortlex-smaller than the
      candidate $(\texttt{s},\texttt{k})$ at $(2,(0,1),0)$, so
      discarding the newcomer loses nothing.
\end{itemize}

\emph{Step 3: pop $(\texttt{k})$} (key $(1,(1),0)$). Accept
$(0,0) \mapsto (\texttt{k})$. Extend:
\begin{itemize}
\item by \texttt{s}: denotation $s \circ k = (s(0), s(0)) = (1,1)$
      --- new; push word $(\texttt{k}, \texttt{s})$ with key
      $(2, (1,0), 0)$.
\item by \texttt{k}: denotation $k \circ k = (0,0)$ --- already
      represented; skip.
\end{itemize}

\emph{Step 4: pop $(\texttt{k}, \texttt{s})$.} Accept
$(1,1) \mapsto (\texttt{k}, \texttt{s})$. Extend:
\begin{itemize}
\item by \texttt{s}: $s \circ (1,1) = (0,0)$ --- represented; skip.
\item by \texttt{k}: $k \circ (1,1) = (0,0)$ --- represented; skip.
\end{itemize}
Frontier empty; the closure completes.

\subsection{The completed table and its evolution}

\begin{center}
\begin{tabular}{@{}cllll@{}}
\toprule
After step & $(0,1)$ & $(1,0)$ & $(0,0)$ & $(1,1)$ \\
\midrule
seeding & --- & --- & --- & --- \\
1 & $\varepsilon_B$ & --- & --- & --- \\
2 & $\varepsilon_B$ & $(\texttt{s})$ & --- & --- \\
3 & $\varepsilon_B$ & $(\texttt{s})$ & $(\texttt{k})$ & --- \\
4 & $\varepsilon_B$ & $(\texttt{s})$ & $(\texttt{k})$ & $(\texttt{k},\texttt{s})$ \\
\bottomrule
\end{tabular}
\end{center}

The final catalog has $\texttt{denotation\_count} = 4$, and the
\texttt{normal\_forms} listing, sorted by shortlex key, is
$\varepsilon_B$, $(\texttt{s})$, $(\texttt{k})$,
$(\texttt{k},\texttt{s})$.

\subsection{Verifying leastness by enumeration}

Because the instance is tiny we can confirm shortlex-leastness against
brute force. The words of length $\le 2$ and their denotations:

\begin{center}
\begin{tabular}{@{}llll@{}}
\toprule
Word & Key & Denotes & Least for its denotation? \\
\midrule
$\varepsilon_B$ & $(0,(),0)$ & $(0,1)$ & yes \\
$(\texttt{s})$ & $(1,(0),0)$ & $(1,0)$ & yes \\
$(\texttt{k})$ & $(1,(1),0)$ & $(0,0)$ & yes \\
$(\texttt{s},\texttt{s})$ & $(2,(0,0),0)$ & $(0,1)$ & no: $\varepsilon_B$ is shorter \\
$(\texttt{s},\texttt{k})$ & $(2,(0,1),0)$ & $(0,0)$ & no: $(\texttt{k})$ is shorter \\
$(\texttt{k},\texttt{s})$ & $(2,(1,0),0)$ & $(1,1)$ & yes \\
$(\texttt{k},\texttt{k})$ & $(2,(1,1),0)$ & $(0,0)$ & no \\
\bottomrule
\end{tabular}
\end{center}

Every denotation's least word is exactly the table's representative.
In particular $(1,1)$ --- the constant to carrier index $1$ --- is
reachable only at length two, and among length-two words only
$(\texttt{k},\texttt{s})$ denotes it, so leastness there is by
uniqueness. Queries now behave as Chapter~9 describes:
$\texttt{normalize}((\texttt{s},\texttt{k})) = (\texttt{k})$,
$\texttt{normalize}((\texttt{s},\texttt{s})) = \varepsilon_B$, and
\texttt{representative\_for} on the arrow $(1,1)$ returns
$(\texttt{k},\texttt{s})$, while on any arrow over an object other
than $B$ it raises \texttt{UnreachableDenotationError}.

\subsection{Order-dependence and the resource guard, exhibited}

Two boundary behaviors of Chapter~8 are visible even at this scale.
First, \emph{representatives depend on the declared generator order}:
re-declaring the catalog with \texttt{k} at rank $0$ and \texttt{s} at
rank $1$ changes no denotation and no reachability, but the length-one
words swap keys, and the constant $(1,1)$ keeps the same spelling with
different key $(2,(0,1),0)$ --- while in a richer catalog a
reordering can change which of two equal-length words wins. Canonical
forms are relative to the versioned catalog identity, exactly as the
claim boundary states. Second, \emph{the guard fails closed}: rerunning
this catalog with \texttt{max\_arrows} $= 3$ proceeds through steps
1--3 and then, at step 4, finds the accepted count already at the cap
\emph{before} inserting $(1,1)$ and raises
\texttt{ClosureLimitExceeded} out of the constructor --- so no catalog
object exists, and no caller can mistake the three-entry table for a
completed closure. The reader who traces these four steps has, in
miniature, executed every branch of the closure engine: acceptance,
both skip paths, the push-side deduplication, completion, and (in the
rerun) the fail-closed guard.
